%!TEX TS-program = pdflatex
\documentclass[11pt,a4paper]{article}

% Packages
\usepackage[margin=1in]{geometry}
\usepackage[hidelinks]{hyperref}
\usepackage{xcolor}
\usepackage{titlesec}
\usepackage{array}
\usepackage{enumitem}

% Colors
\definecolor{darkblue}{HTML}{003366}
\definecolor{lightgray}{HTML}{808080}

% Section formatting
\titleformat{\section}
  {\normalfont\Large\bfseries\color{darkblue}}
  {}
  {0em}
  {}
  [\titlerule\vspace{0.3cm}]

\titleformat{\subsection}
  {\normalfont\normalsize\bfseries}
  {}
  {0em}
  {}
  [\vspace{0.1cm}]

% Remove page numbers
\pagestyle{empty}

% Custom commands
\newcommand{\entry}[4]{
  \noindent\textbf{#1} \hfill \textit{#2} \\
  \noindent\textit{#3} \\
  \noindent #4 \\[0.3cm]
}

\newcommand{\simpleentry}[2]{
  \noindent\textbf{#1} \hfill \textit{#2} \\[0.2cm]
}

\newcommand{\activity}[2]{
  \noindent$\bullet$ #1 \textit{#2} \\[0.15cm]
}

\newcommand{\award}[1]{
  \noindent$\bullet$ #1 \\[0.15cm]
}

\begin{document}

% Header
{\centering
\Huge\textbf{Giang Vu Le} \\
\vspace{0.2cm}
{\large PhD Student in Economics} \\
\vspace{0.2cm}
\normalsize
IMT School for Advanced Studies Lucca, Italy \\
Email: \href{mailto:giangvu.le@imtlucca.it}{giangvu.le@imtlucca.it} \\
Website: \href{https://legiangvu.github.io}{legiangvu.github.io} \\
\vspace{0.3cm}
}
\noindent\textit{Last Updated: \today} \\
\vspace{0.3cm}

% Education
\section*{EDUCATION}

\entry{PhD in Economics}{2024 -- Present}
{IMT School for Advanced Studies Lucca, Italy}
{}

\entry{Master of Arts in Economic Policy}{2022 -- 2024}
{Central European University, Vienna, Austria}
{}

\entry{Bachelor of Engineering in Aeronautical Engineering}{2018 -- 2022}
{Ho Chi Minh City University of Technology, Vietnam}
{}

% Research & Publications
\section*{RESEARCH}

\subsection*{Work in Progress}

\noindent\textbf{``Picking the Right Winners? Optimal Policy Learning and Subsidies''} \\
with Armando Rungi. Draft in preparation. \\
\noindent\textit{Abstract.} We propose applying an optimal policy learning framework to identify \textit{ex ante} allocation rules of state aid using firm-level financial accounts. Our compelling case study is the tendering of state aid in Germany for the energy transition. First, we evaluate the factual policy, and then we compare it with two alternatives derived from distinct social preferences: a maximum-employment and a maximum-growth policy. We apply a Doubly Robust learner to estimate average and conditional treatment effects (CATEs), and then study different scenarios with constrained and unconstrained policies. We find that the factual policy is always suboptimal, while it shows only an average increase in capital stock and a decrease in productivity. When we examine the distribution of CATEs, we observe substantial heterogeneity, with both negative and positive effects on relevant outcomes. Once we consider alternative optimal policies, a trade-off emerges between pursuing maximum employment and maximum growth. At some point, the pursuit of maximum employment begins to destroy economic value. Also, we find that the maximum-employment policy would disproportionately target larger firms, while the value-added-maximizing policy would favor smaller firms, leading to higher productivity. Finally, we discuss how moving away from discretionary or size-based selection toward objective, data-driven allocation rules can significantly improve the efficiency and impact of public policies.

% % Skills
% \section*{SKILLS}

% \noindent\textbf{Programming Languages:} Python, R, Stata, MATLAB \\
% \noindent\textbf{Languages:} English (Fluent), Vietnamese (Native), French (Intermediate), Italian (Beginner)


\section*{PROFESSIONAL ACTIVITIES}

\activity{Organizer of the CIML Training \& Conference ``Frontiers of Causal Inference and Machine Learning''}{2026}

\activity{Organizer of the IMT Lucca Reading Group - Policy challenges: international economics, geopolitics, and global innovation}{2025 -- 2026}

\activity{Organizer of the Global Firms and Knowledge Diffusion Workshop}{2025}

\activity{Organizer of the CIML Training \& Conference ``Frontiers of Causal Inference and Machine Learning''}{2025}

% Professional Activity
\section*{PRESENTATIONS \& CONFERENCES}

\activity{CIML Training \& Conference ``Frontiers of Causal Inference and Machine Learning'', Lucca, Italy}{2025}

% \vspace{0.3cm}
% \noindent\textbf{Professional Memberships} \\
% \noindent To be updated with memberships

% Awards
\section*{AWARDS \& HONORS}

\award{Outstanding Thesis Award in the MA in Economic Policy in Global Markets Program (2023/24), Central European University, Austria}

\award{Outstanding Academic Achievement Award for the highest GPA in the department (2023/24), Central European University, Austria}

\end{document}
